\documentclass[a4paper]{article}
\usepackage{amsfonts}
\usepackage{amsmath}
\usepackage{amssymb}
\usepackage{graphicx}     

\begin{document}
  
\noindent \large Auteur: Kyan Van den Eynde \\
\noindent \large Studierichting: Informatica\\
\noindent \large Oefeningenreeks 2A nr. 4\\

\medskip

\normalsize

Als $m=0$ dan gaat het over een constante functie en geldt voor iedere $a\in\mathbb{R}$ en voor iedere $\varepsilon > 0$ dat:\\

$|f(x)-f(a)| = |q-q| = 0 < \varepsilon$\\

Als $m\neq 0$ dan:\\

Kies $a \in \mathbb{R}$, Pak $\varepsilon > 0$, Stel $\delta = \dfrac{\varepsilon}{|m|}$\\

$|f(x)-f(a)| = |mx+q-ma-q| = |m(x-a)| = |m|\cdot |x-a| < |m| \cdot \dfrac{\varepsilon}{|m|} = \varepsilon$\\

$|x-a| < \delta$ dus $|m(x-a)| < \varepsilon$\\

\begin{thebibliography}{999}
%\bibitem{a} Referentie
\end{thebibliography}
\end{document} 
